\documentclass[11pt]{article}
\usepackage[margin=1in]{geometry}
\usepackage{tgtermes}
\usepackage{setspace}
\setstretch{1.15}
\usepackage{microtype}
\usepackage{amsmath,amssymb}
\usepackage[hidelinks,pdftitle={What an intervention test would have to show},pdfauthor={}]{hyperref}
\usepackage[authoryear,round]{natbib}
\setcitestyle{aysep={}}
\pagestyle{plain}
\begin{document}
\begin{center}
{\Large What an intervention test would have to show\par}
\vspace{0.7em}
\end{center}

\begin{abstract}
A process-function model predicts that every replacement of the local maps has exactly one history, and a companion theorem says that the same condition is necessary and sufficient for a canonical closed loop to stay unitary under local unitaries and matched finite ancillas. Those statements fix the pattern an intervention test of the model would have to recover, if a local map has been identified with a manipulation and a history with a record. They are not themselves the test. Evidential warrant for the model as a description of a target would be a match, under those identifications, between the predicted pattern and the records. Calculations internal to the model, including a negative value of a linear functional on a known board in one basis, check consequences of definitions. They do not supply the identifications, and they do not confirm or refute a claim about a device. A finite search shows that an explicit table can be checked. It is a precondition for using the table in a test, not a substitute for one. The mathematics is an alternative adaptation of results already public. No new theorem is offered, and no target is claimed to have been tested.
\end{abstract}

\noindent\textit{Keywords:} intervention; models; evidential warrant; process functions; robustness

\section{The gap between a prediction and a test}
\label{sec:gap}

The finite mathematics used here is already public. A map on finite alphabets is a process function when every tuple of local functions has exactly one fixed point, called a history. A canonical extension of any such map, process function or not, is a basis permutation of a source-sink space. The closed loop built from that extension is unitary for every local unitary and every matched finite ancilla if and only if the map is a process function. A known switchboard induces, in the computational basis, a distribution on which a stated linear functional is negative, while the same functional is nonnegative on three-bit process functions in an enumeration that is not repeated in this file. None of this is re-proved below. The question is evidential. What would it take for these statements to be tested against a target, and which of the existing calculations already are such a test?

The short answer is that none of them are. They determine a pattern of outcomes. A test would have to realize the interventions the pattern quantifies over, record the outcomes the pattern talks about, and compare. \citet{kuorikoski2007} argue, in a different science, that robustness of a modeling result can say which part of a model a later empirical test would be testing, and that it does not replace the test. They quote Orzack and Sober's claim that the robustness of a theorem is not by itself a reason to believe it, and they report Sugden's claim that comparisons among models do not license an inductive leap to the world. I use that limit, from the preprint I read. I do not re-read Orzack and Sober or Sugden, and I do not treat Kuorikoski, Lehtinen, and Marchionni as having discussed process functions. They had not.

A companion adaptation of the same mathematics, written for a different journal, uses Karaca's distinction between the robustness of a procedure and the robustness of a system that would carry the procedure out \citep{karaca2022}. That distinction is assumed here and not re-argued. The present question is narrower and prior to implementation. Even before one asks whether a system is adequate, one can ask what an intervention test of the model would consist in, and whether anything already written down is evidence that the model holds of a target. The answer to the second question is no. The answer to the first is a short list of identifications. Stating the list is the contribution of this adaptation. It is not a new theorem and not an independent study.

\section{The pattern the model predicts}
\label{sec:pattern}

Fix finite alphabets. A wiring \(w\) sends a tuple of outputs to a tuple of inputs, one component per party. A local function at party \(k\) sends that party's input alphabet to its output alphabet. The wiring is a process function when, for every choice of local functions, the map that sends an output tuple \(o\) to the tuple of local-function values at \(w(o)\) has exactly one fixed point \citep{baumeler2016}. The quantifier is universal. A single solved instance is not the prediction.

Three small cases show what the prediction says, and they are calculations, not episodes. If one party copies a bit and the other flips it, there is no output pair that satisfies both local maps. The model predicts that a test which had truly set those two maps would find no history. If the second party copies rather than flips, both bit values are histories. The model predicts two records, not one. If the wiring is the three-bit rule \(w(a,b,c)=(\lnot b\land c,\ \lnot c\land a,\ \lnot a\land b)\) and each party copies, the only history is \((0,0,0)\), and the same uniqueness holds for every other substitution of the binary functions. The model predicts one record for each of those sixty-four substitutions. A test that found two records for one of them would refute the claim that this wiring is a process function, under the identifications of Section~\ref{sec:id}. A test that found one record for one substitution would not establish the universal claim.

The unitary theorem is a second prediction, about a larger class of inputs. The canonical operator \(U\lvert o,e\rangle=\lvert w(o)+e,o\rangle\) is a permutation for every wiring. Inserting local unitaries on matched finite ancillas yields a closed loop that is unitary for every such insertion if and only if the wiring is a process function. The prediction for a test, if the local unitaries have been identified with manipulations and the loop's unitarity with a recorded reversibility, is that the recorded reversibility holds for all of those manipulations exactly when the classical test would have found one history in every case. The prediction is not that some reversible encoding exists. The encoding exists for the inconsistent wirings as well.

Two limits belong in the prediction itself, so that a test is not asked to confirm something the model does not say. The argument uses finite dimension. It does not say that every exclusive self-map of an infinite product of bits has a fixed point. That question, for Baire-class-one maps, is open and is not part of the prediction. The converse direction of the unitary theorem uses one family of local unitaries, the reversible encodings of functions. It does not say that an arbitrary purification is unitary only for process functions. A test designed around a different dilation would be a test of a different claim.

\section{The identifications a test would need}
\label{sec:id}

An intervention test, as the phrase is used here, is not a formal intervention in the sense of a theory of causal graphs. It is a comparison with the following structure.

First, some manipulation available in the target is identified with a local function, or with a local unitary, in the model. The identification has to say which manipulation corresponds to which map, and it has to say that the manipulation sets that map without silently setting the others. The process-function quantifier treats the local maps as independently replaceable. If in the target one cannot set one party's map without setting another's, the quantifier is not the quantifier of that target, and a failure of unique solvability under independent replacement is not a failure of the target. This is a precondition of the test, not a result of it. Directed modularity, in the sense used by a different adaptation of this mathematics, is a different independence condition: untouched equations stay invariant when an arrow is broken \citep{hausman1999}. I do not re-derive that paper. I use only the contrast. The process-function model assumes a different freedom, freedom to replace local maps, including replacements that leave cyclic dependence in place. A test of this model has to justify that freedom for the target. A test of directed modularity would justify a different freedom. Passing one does not pass the other.

Second, some record available in the target is identified with a history. The history is a tuple of outputs that is a fixed point of the composed maps. A record that is not a candidate for that tuple, or that mixes several tuples, is not yet the outcome variable of the test. The identification can be wrong even when the arithmetic is right. Finding the number \((0,0,0)\) in a log is not finding the history of the three-bit rule unless the log's coordinates are the rule's outputs.

Third, the test applies the manipulations in the class the prediction quantifies over, or in a stated subclass, and compares the records with the predicted pattern. For the universal claim, the class is every local replacement. A test of a subclass supports at most the restricted claim. Sixty-four substitutions of binary functions on three bits is a finite class. It can be run, if the identifications are in place. It cannot be replaced by running one substitution.

Fourth, the comparison has a stated tolerance for what counts as a miss. In the exact model, a second fixed point is a miss, and a missing fixed point is a miss. A laboratory record will have noise. The model as stated has no noise term. Extending it by a noise term would be a different model, and the warrant would then be warrant for that extension. This paper does not supply the extension. It records that the exact prediction cannot be compared with noisy records until someone says how much disagreement the exact claim tolerates. That saying is part of the test's design.

These four items are individually familiar. Together they are what is missing from a file that contains only the theorems. The theorems are the content of the third item's predicted pattern. They are not the first, second, or fourth item.

\section{What the existing calculations are evidence for}
\label{sec:calc}

Inside the model, the theorems are derivations. They are evidence that the stated pattern follows from the stated definitions, for a reader who checks the proof. They are not evidence that a target instantiates the definitions. \citet{kuorikoski2007} put the corresponding point about economic models by agreeing that robustness is not, by itself, empirical confirmation, and by insisting that it can still show which assumptions a result does not depend on. The process-function theorem shows that one history per replacement does not depend on which replacement is chosen, when the universal claim holds, and that unitarity of the canonical loop stands or falls with that claim. A later test could be aimed at that stable biconditional rather than at a single gate setting. The theorem does not carry out the test. This is the limited role Kuorikoski, Lehtinen, and Marchionni allow, transferred to a model they did not discuss. If the transfer fails, the failure is that their defense was about the security of inferences from uncertain economic assumptions, not about a biconditional proved from a finite definition. I do not need their positive defense. I need the negative point they share with the authors they quote: a result's stability inside a model is not a reason, by itself, to believe a claim about a target.

The bare permutation is easy to misread as evidence of realizability. It can be written down for every wiring, including the two-bit wirings that have no history or two histories. Its existence is a theorem of linear algebra on a finite set. A test of realizability would ask whether a device's transitions are those of the permutation, or whether its records match the histories. The permutation's existence does not answer either question. It is evidence only that the relabeling is a well-defined operator.

The finite search is evidence of a different internal fact. For an explicit table, the set of output tuples is finite, and a machine can count fixed points for each local replacement. The search shows that the prediction is decidable from the table. Decidability matters for test design: one cannot compare records with a prediction one cannot compute, and here one can compute it. Decidability is not the comparison. A search that prints ``one fixed point'' has checked the model against itself. It has not checked a device. It also does not show that the same decision procedure is fast for a family of tables the search was not given, and it does not show that the wiring computes a function the search does not compute. Those would be further claims. They are not supported by the existence of the search.

The known-board calculation is the case most likely to be called empirical, because it produces probabilities. A table of exponents defines a unitary for each output string. A source vector is stated. The conditional distribution is the squared amplitude in the computational basis. A linear functional of that distribution equals \(1/4-(3/16)\sqrt{2}\), which is negative. If an external enumeration is granted, the functional is nonnegative on every three-bit process function, so the distribution is not a mixture of them. The calculation uses the Born rule as a definition of the distribution, not as a summary of frequencies from a sample. No device was run. The negativity is evidence that this particular mathematical distribution lies outside the hull, conditional on the enumeration. It is not evidence that a laboratory frequency would lie outside the hull, and it is not evidence that the hull fails as a description of a target. A test of that further claim would have to identify the board's unitaries with manipulations and the computational-basis outcomes with records, and then sample. The identification is not in the table. The sample is not in the file.

Purity of the board, in the public note, is a finite cancellation of a stated number of cyclotomic coefficients, or an equivalent multilinear check. That cancellation is not reproduced here. Even if it is granted, it is evidence that the board defines a pure process in the model's sense. Purity in that sense is unitarity of a loop under local unitaries. It is not a laboratory report that a state was pure.

\section{How the model could be disconfirmed}
\label{sec:dis}

It helps to say what a miss would look like, so that the claim is not insulated. Suppose the identifications of Section~\ref{sec:id} are granted for the three-bit rule, and suppose a manipulation corresponding to one substitution of local functions yields two distinct records that both satisfy the fixed-point equation, or yields none. Then the universal claim is false of that target. The arithmetic of the rule is not thereby false. The rule, as a mathematical map, still has one fixed point for that substitution. What fails is the claim that the target's records are the rule's histories under the stated identification. That is ordinary falsification of an empirical application. It is available only after the identification is fixed. Without the identification, two records are not a counterexample to the rule. They are two numbers.

A miss for the unitary prediction is different and easier to confuse with a miss for the classical one. Suppose a device is claimed to realize the canonical loop, and a local unitary insertion yields a map that is not unitary, while the classical substitutions still have unique histories. Then the biconditional's application to that device fails. The theorem is not refuted, because the theorem is about the operator defined from \(w\), not about the device. The refutation, if the identifications hold, is of the claim that the device's loop is that operator. Conversely, a device could have unique classical histories and fail to be the canonical unitary construction, because some other dilation was built. The theorem says the canonical one is unitary when the histories are unique. It does not say that every unitary device with unique histories is this construction.

A miss for the witness claim would be a demonstration that the printed functional is nonnegative, or that the distribution was miscomputed from the table. That miss is internal. It can be checked by arithmetic. It does not require a laboratory, and a laboratory cannot repair an arithmetic error in the table. Keeping these misses apart is the point of the section. An empirical miss, an implementation miss, and an arithmetic miss are three different events. A paper that calls all three ``the model failed'' has not yet said what was tested.

\section{A worked protocol, still not a test}
\label{sec:protocol}

The four identifications can be written as a protocol for the three-bit rule without thereby running it. The protocol is useful because it shows where the public mathematics stops and the evidential work would start. It is also a check on length of claim. A protocol that cannot be stated has not been shown to be a possible test. A protocol that can be stated, and has not been carried out, is an untested application.

Take the wiring \(w(a,b,c)=(\lnot b\land c,\ \lnot c\land a,\ \lnot a\land b)\). There are four functions from a bit to a bit. A full classical protocol lists the sixty-four triples of functions. For each triple it requires a manipulation that sets party \(k\)'s input-output behavior to the chosen function, independently of the other parties' settings, and a record of the three output bits after the setting. The prediction, already proved for the rule as a mathematical object, is that exactly one triple of output bits satisfies the fixed-point equation, and that the protocol's record is that triple. The protocol's empirical content is exhausted by whether the record matches. Repeating the proof inside the protocol does not add a datum.

A smaller protocol is available and supports less. One may test only the three copying functions, and predict the record \((0,0,0)\). A match supports the instance. It does not support the universal claim, because the other sixty-three substitutions were not set. One may also test the two-bit examples as negative controls inside the same laboratory, if the identifications extend to them. The copy-and-flip pair should yield no record that satisfies both maps. The copy-and-copy pair should yield two. A laboratory that found a unique record for copy-and-flip would have a miss of the sort in Section~\ref{sec:dis}, and the miss would concern the identification or the device, not the theorem that the pair of maps has no fixed point. The negative controls matter because they stop a weak form of confirmation: a device that prints \((0,0,0)\) for every setting is not confirming the three-bit rule. It is failing the controls. Orzack and Sober's remark, as quoted by \citet{kuorikoski2007}, that robustness is not by itself a reason to believe the robust claim, has a practical version here. A record that agrees with the robust case and has not been compared with the non-robust cases has not yet been shown to track the feature that distinguishes them.

The unitary protocol is stricter and easier to fake in prose. It would require manipulations corresponding to local unitaries, not only to functions, on ancillas whose dimensions match the alphabets, and a record of whether the loop map is unitary. Unitarity is not a single click. It is a relation among amplitudes. A protocol that records only computational-basis frequencies has not recorded unitarity. It has recorded a classical marginal. The known-board distribution is such a marginal, computed rather than sampled. Using it as if it were the record of the unitary protocol would confirm a different, weaker claim, and only if the frequencies were sampled from a device rather than calculated from a table. The biconditional says the loop is unitary for all local unitaries exactly when every functional substitution has one history. A protocol that checks one unitary and one substitution has checked neither side of that biconditional.

Two further design choices belong in the protocol because they change what is warranted. The first is the order of settings. The model has no temporal order among the local maps. A laboratory procedure has an order, because manipulations take time. If the record depends on that order, the target has structure the model does not state. Agreement in one order and disagreement in another would be a miss relative to the claim that the model is the whole story, and it would not be a miss relative to the weaker claim that the model describes the unordered fixed points when the order is held fixed. The protocol has to say which claim is under test. The second choice is the source of the local maps. If the maps are set by the experimenter, the quantifier ranges over interventions. If the maps are merely observed to vary, the variation may be entangled with the outputs, and the fixed-point equation is not an intervention test. It is a joint observation. The process-function definition does not by itself say that observed variation is an intervention. A protocol that treats every observed change as a setting has added a premise. The premise needs its own warrant. The theorem does not contain it.

Nothing in this protocol is an instruction to build a device, and nothing in it asserts that the three-bit rule is a good model of any laboratory that exists. The protocol's role is to make the evidential conditional explicit. If the identifications are granted and the records match across the stated class, the application is supported to that extent. If they do not match, the application is not. The public files establish the consequent of the prediction under the mathematics. They leave the antecedent, and the comparison, undone.

\section{Objections}
\label{sec:obj}

The first objection says that the identifications are pedantic because the model was built to describe interventions. A model can be intended for interventions and still not have been tested. Intention fixes what the identifications ought to look like. It does not establish that they hold for any particular device. The three-bit rule has been discussed as a logical possibility of a cyclic dependence with unique solutions. Discussion is not an intervention. Until someone sets the local maps and writes down the records, the evidential status of the rule as a description of a target is the status of an untested model. That status is not a complaint about the rule's interest. It is a description of what has been done.

The second objection says that the universal quantifier is itself a kind of evidence, because a claim that survives every intervention is better supported than a claim that survives one. The objection borrows the warrant of an empirical robustness argument and spends it on a quantifier inside a definition. If every intervention in a class has been performed and the outcome was as predicted, the universal claim is supported to the extent the class was the relevant one and the identifications held. If every intervention has been quantified over in a proof, the proof supports the conditional claim that if the identifications hold, the pattern is such and such. The proof does not perform the interventions. \citet{karaca2022} distinguishes the invariance of a procedure's function under varying inputs from the convergence of detection results. The quantifier is the first kind of invariance, for a procedure that exists on paper. It is not the second.

The third objection says that refusing to count the witness as evidence makes the model unfalsifiable, which is the charge Orzack and Sober press against robustness. The charge would stick if no pattern of records could count against an application. Section~\ref{sec:dis} states patterns that would. What is refused is the claim that a calculation already is those records. Falsifiability of an application requires a bridge to records. It does not require that the bridge have been crossed. A model with a stated bridge and no data yet is untested. A model with no conceivable bridge is a different problem. This model has a conceivable bridge. It has not been crossed in the public drafts.

The fourth objection says that publishing the table of exponents is sharing data, and shared data are evidence. The table is a specification of a unitary family. Sharing it lets another person repeat the contraction. Repeatability of arithmetic is a virtue of a proof. It is not a sample size. The journal's data practices, discussed in the audit of this draft rather than settled here, ask authors to say where calculations live. They do not convert a table into an experiment.

The fifth objection says that this paper repeats the robustness paper and the explanation paper under a new label. The overlap is real and should be owned. All three use one biconditional. The explanation paper asks when an admissibility constraint would explain a target phenomenon, in a sense taken from Bokulich, and answers that a justificatory step is still missing. The robustness paper asks which robustness notion the biconditional is, and answers that it is procedure-level invariance, not implementation. This paper asks what an intervention test would have to contain, and answers with four identifications and a separation among three kinds of miss. A reader who thinks the third question is only a restatement of the first two is right that no new mathematics has been introduced, and right that the papers are alternative adaptations. The reader is wrong if the objection means that a list of identifications required for a test is the same claim as a taxonomy of robustness or a set of conditions on explanation. The list can be rejected while the taxonomy is kept, and conversely. That is as much distinctness as the shared mathematics will support. It is not advertised as more.

\section{Instance warrant is not universal warrant}
\label{sec:instance}

Even after a protocol is run, the warrant has a scope, and the scope is easy to inflate. A match on one substitution warrants the instance, under the identifications and the tolerance. The model’s characteristic claim is universal. Support for a universal claim from a proper subclass is support only if one has a reason to think the subclass is representative. The mathematics provides no such reason. The interesting substitutions are not a random sample from the sixty-four. The copying substitution is one point. The bit-flip failures are other points. A laboratory that confirms the copying case and the two negative controls has more warrant than a laboratory that confirms only the copying case, and it still has not confirmed the remaining substitutions. There is no theorem in the public note that says a few instances imply the rest for an unknown target. The implication holds inside the model only because the wiring is fully specified and the proof covers every case. A target is not specified in that way until the identifications are in place for each case one wishes to claim.

The same inflation threatens the unitary side. One local Hadamard gate is a natural first manipulation, because it is not diagonal in the computational basis and therefore is not a classical function in disguise. A loop that remains well behaved under that one gate is an instance. The theorem’s quantifier is every local unitary. One gate does not stand in for the quantifier. Nor does the computational-basis witness. That witness is a particular source vector and a particular family of unitaries, read in one basis. Its interest, inside the model, is that the induced distribution can leave the convex hull. Its interest is not that it samples the space of local unitaries. Treating it as a representative intervention would be a sampling claim the calculation does not make.

Scope also runs in the other direction. A failed instance, under sound identifications, is enough to reject the universal application. One does not need to run the other substitutions to retract the claim that every substitution matches. This asymmetry is ordinary and should be kept. It is why the negative controls in Section~\ref{sec:protocol} are efficient, and why a single unmatched record is informative. It is not a reason to treat a single matched record as if it had the same power. Confirmation of a universal claim and refutation of a universal claim are not symmetric in the number of instances they require. The public proofs are symmetric biconditionals about maps. The evidential relation between instances and the universal application is not that biconditional. Confusing the two is how a checked example gets described as a tested theory.

\section{Conclusion}
\label{sec:end}

The process-function condition and the canonical unitary equivalence are predictions about patterns. An intervention test of those predictions, as descriptions of a target, would identify manipulations with local maps, records with histories, a tolerance with the exact equality the model uses, and then compare. The public derivations do the last of these only inside the model, where the ``records'' are fixed points computed from the maps. They show what the comparison should look like. They do not perform it on a target.

Evidential warrant for an application would be the outcome of that comparison. Warrant for the theorems is a different thing and is not in dispute here: a reader checks a finite proof. The two warrants are not interchangeable. A negative functional on a known board is warrant, at most, for a mathematical separation, and only together with an enumeration this file does not repeat. A search that decides a finite table is warrant that the prediction can be computed. A permutation that exists for inconsistent wirings is warrant that reversibility of a relabeling is not the process-function condition. None of these is warrant that a device has been tested, explained, or built.

What would increase warrant for an application is dull and specific. State the identifications. Run the substitutions, or a declared subclass. Say what disagreement counts as a miss. Until that is done, the correct evidential description of the public mathematics is that it is available for a test and has not been put to one.

\section*{Use of a language model}
Grok Bot, a Cursor cloud agent, drafted this manuscript, including the organization of the argument, on 8 November 2026. The use went beyond copy editing. The model version is not recorded and is not invented. The human author remains responsible for the text. Verification is pending. The agent is not an author. A complete verbatim log of prompts is not retained. The operative instruction was to write a focused argument, for this journal, about modeling, intervention tests, and evidential warrant, as an alternative adaptation of the same finite results, without a new theorem.

\begin{thebibliography}{99}
\bibitem[Baumeler and Wolf(2016)]{baumeler2016} Baumeler, \"A., \& S. Wolf. (2016). The space of logically consistent classical processes without causal order. \textit{New Journal of Physics, 18}, 013036. \url{https://doi.org/10.1088/1367-2630/18/1/013036}
\bibitem[Hausman and Woodward(1999)]{hausman1999} Hausman, D. M., \& J. Woodward. (1999). Independence, invariance and the causal Markov condition. \textit{British Journal for the Philosophy of Science, 50}, 521--583. \url{https://doi.org/10.1093/bjps/50.4.521}
\bibitem[Karaca(2022)]{karaca2022} Karaca, K. (2022). Two senses of experimental robustness: Result robustness and procedure robustness. \textit{British Journal for the Philosophy of Science, 73}, 279--298. \url{https://doi.org/10.1093/bjps/axy031}
\bibitem[Kuorikoski, Lehtinen, and Marchionni(2007)]{kuorikoski2007} Kuorikoski, J., Lehtinen, A., \& Marchionni, C. (2007). \textit{Economics as robustness analysis}. PhilSci-Archive. \url{https://philsci-archive.pitt.edu/3550/}
\end{thebibliography}
\end{document}
